\documentclass[12pt]{article}
\usepackage[T1]{fontenc}
\usepackage[utf8]{inputenc}
\usepackage{lmodern}
\usepackage{textcomp}
\usepackage{enumitem}
\usepackage{geometry}
\geometry{margin=1in}
\begin{document}
\begin{center}
Answer Key - Astronomy - Graduate Level Assessment 1 \
\small (Answers are ordered exactly as in the question paper.)
\end{center}

\section*{Multiple Choice}
\begin{enumerate}[label=\arabic*.]
\item \textbf{Answer:} B
\\ \textit{Options:} A) To measure interstellar mass.; B) To measure galactic distances.; C) To measure galactic energy-density.; D) To measure interstellar density.
\\ \textit{Note:} Cepheid stars are relevant for astronomers because their predictable periodic brightness variations allow for accurate distance measurements within galaxies, serving as a crucial tool for establishing the scale of the universe.
\item \textbf{Answer:} D
\\ \textit{Options:} A) a dark ring visible from Earth composed of dark dusty particles; B) the most opaque ring of Saturn made of highly reflective ice particles; C) the widest ring of Saturn located between two large ring gaps; D) a large gap visible from Earth produced by an orbital resonance with the moon Mimas
\\ \textit{Note:} The Cassini division is the prominent gap in Saturn's rings that arises due to gravitational interactions between the particles in the rings and the moon Mimas, resulting in an orbital resonance that clears out a region between the A and B rings.
\item \textbf{Answer:} D
\\ \textit{Options:} A) the Sun mainly emits blue light.; B) the atmosphere absorbs mostly blue light.; C) molecules scatter red light more effectively than blue light.; D) molecules scatter blue light more effectively than red light.
\\ \textit{Note:} The correct answer is based on Rayleigh scattering, which demonstrates that shorter wavelengths of light, such as blue, are scattered more by atmospheric molecules than longer wavelengths like red, resulting in the blue appearance of the sky.
\item \textbf{Answer:} A
\\ \textit{Options:} A) Distance at which one astronomical unit measures one arcsecond from Earth.; B) Orbital distance of the solar system around the center of the Milky Way in one year.; C) Effective distance of the solar wind (i.e. the radius of the heliosphere).; D) Historical distance to the brightest star Sirius.
\\ \textit{Note:} One parsec is defined as the distance at which one astronomical unit subtends an angle of one arcsecond as viewed from Earth, establishing a direct relationship between distance and angular measurement in the context of astronomical observations.
\item \textbf{Answer:} A
\\ \textit{Options:} A) 750 million years.; B) 2 billion years.; C) 3 billion years.; D) 4.5 billion years.
\\ \textit{Note:} The surface of Venus is approximately 750 million years old, based on geological studies that suggest it has been resurfaced by volcanic activity relatively recently in the planet's history, indicating a younger surface compared to other celestial bodies in the solar system.
\end{enumerate}

\section*{True / False}
\begin{enumerate}[label=\arabic*.]
\setcounter{enumi}{5}
\item \textbf{Answer:} False
\\ \textit{Options:} A) True; B) False
\\ \textit{Note:} Asteroids do not emit their own radiation; instead, they reflect sunlight, which is the primary reason for their visibility in the night sky.
\item \textbf{Answer:} False
\\ \textit{Options:} A) True; B) False
\\ \textit{Note:} Ganymede is the largest moon in the solar system, but it is not the closest to Jupiter; that title belongs to Io, which orbits at a shorter distance from the planet.
\item \textbf{Answer:} False
\\ \textit{Options:} A) True; B) False
\\ \textit{Note:} The statement is false because weight is directly proportional to gravitational acceleration, so on a planet with half the gravity of Earth, a 5 kg sand bag would weigh only 25 N, not 50 N.
\item \textbf{Answer:} True
\\ \textit{Options:} A) True; B) False
\\ \textit{Note:} Fresh impact craters on solid surfaces typically exhibit ejecta and raised rims due to the explosive nature of the impact, which displaces material outward and upward, creating these distinctive features around the crater.
\item \textbf{Answer:} True
\\ \textit{Options:} A) True; B) False
\\ \textit{Note:} Dione is a moon of Saturn, whereas the Galilean moons are specifically the four largest moons of Jupiter: Ganymede, Callisto, Europa, and Io.
\end{enumerate}

\section*{Long Form Response}
\begin{enumerate}[label=\arabic*.]
\setcounter{enumi}{10}
\item \textbf{Answer:} Pluto's orbital period of approximately 250 years has significant implications for its classification as a dwarf planet and our understanding of the Kuiper Belt. This extended orbital duration highlights the dynamic and diverse nature of trans-Neptunian objects, emphasizing that not all celestial bodies in this region conform to traditional planetary characteristics. Furthermore, Pluto's elongated and inclined orbit, along with its long period, reflects the gravitational influences of nearby objects and the complex history of the Kuiper Belt's formation. As a result, Pluto's classification as a dwarf planet underlines the necessity for a nuanced understanding of planetary criteria, particularly in light of the diverse array of bodies that populate the Kuiper Belt. Thus, Pluto serves as a key case study in the ongoing discourse regarding planetary classification and the evolutionary processes of our solar system's outer regions.
\\ \textit{Note:} correct because it connects Pluto's lengthy orbital period to its unique classification as a dwarf planet and emphasizes the complexity of the Kuiper Belt, illustrating how such characteristics challenge traditional definitions of planets.
\item \textbf{Answer:} Asteroids and comets play a crucial role in the context of the Solar Nebular theory as they represent the remnants of the primordial planetesimals that formed during the early stages of the solar system. These bodies, which never coalesced into full-fledged planets, provide valuable insights into the conditions and materials present in the protoplanetary disk from which the planets ultimately emerged. Their composition, often rich in volatile compounds and organic materials, suggests that they could have contributed to the delivery of water and potentially life's building blocks to terrestrial planets, particularly Earth. Additionally, studying asteroids and comets helps astronomers understand the processes of accretion and differentiation that shaped planetary bodies, making them not only remnants of the past but also key to unraveling the history of our solar system's formation.
\\ \textit{Note:} correct because it accurately describes asteroids and comets as remnants of planetesimals from the early solar system, emphasizing their significance in understanding the composition and conditions of the protoplanetary disk that contributed to planetary formation.
\end{enumerate}

\vfill
\noindent\textit{End of Answer Key}
\end{document}